% !TeX root = ../../Book4.tex
% 纲要标题：### 21.2 完美复形
% 纲要标签：b4:sec:2002
% 纲要位置：计划纲要.md，第 3794 行。

\section{完美复形}
\label{b4:sec:2002}

真正有限的导出对象需要有界有限生成投射模型。厚生成描述它的有限构造，紧性则用它到任意直和的态射来识别同一类对象。前两小节只涉及有界模型与有限构造；第三小节的任意余积与紧对象反向刻画处在无界导出范畴，依赖\secref{b4:sec:1804}建立的半自由替换及\cref{b4:thm:semifree-resolution,b4:prop:k-resolution-comparison}。


% 纲要内容（仅作写作计划，尚非教材正文）：
%
% * 有限投射复形
% * 厚子范畴
% * 紧性概念的适用背景
%

\subsection{有限投射复形}
\label{b4:subsec:2002:01}

一个复形的同调可能仅在有限个次数非零，但计算这些同调所需的投射消解却可能无限延伸。要在导出范畴中表达真正有限的投射数据，应把有限性加在一个可选取的投射模型上。

\begin{definition}\label{b4:def:perfect-complex}
设 \(R\) 为含幺环，\(X\in D(R\text{-}\mathrm{Mod})\)。若 \(X\) 在该导出范畴中同构于某个有界复形 \(P\)，且每个 \(P^n\) 都是有限生成投射左 \(R\)-模，称 \(X\) 为\term[wanmeifuxing]{完美复形}[perfect complex]。全部完美复形组成的全子范畴记为 \(\operatorname{Perf}(R)\)。
\end{definition}
\symidx{\(\operatorname{Perf}(R)\)：环上的完美复形范畴}

\begin{proposition}\label{b4:prop:perfect-homotopy-model}
设 \(R\) 为含幺环，以 \(R\text{-}\mathrm{proj}\) 表示有限生成投射左 \(R\)-模的满子范畴。从同伦范畴到导出范畴的自然函子限制为三角等价
\[
 K^b(R\text{-}\mathrm{proj})\ \simeq\ \operatorname{Perf}(R).
\]
完美复形对有限直和、移位和导出范畴中的映射锥封闭。
\end{proposition}
\begin{proof}
由\cref{b4:prop:resolution-computes-derived-hom}，有界投射复形之间的导出态射就是复形映射的同伦类，所以所给函子全忠实；本质满则是完美的定义。有限直和和移位保持项的有限生成投射性。对完美对象之间的态射，先取有界有限生成投射模型，再用全忠实性把态射表示成复形映射。其映射锥依然有界，每项是两个有限生成投射模的直和，因此仍然完美。好三角来自这样的锥，故等价保持三角结构。
\end{proof}

例如，\(\mathbb Z/m\mathbb Z\) 置于零次是 \(\mathbb Z\) 上的完美复形，因为它有两项有限自由消解 \(\mathbb Z\xrightarrow{m}\mathbb Z\)。这并不要求零次同调本身投射。

\begin{example}\label{b4:exm:bounded-not-perfect}
设 \(k\) 为域，\(R=k[\varepsilon]/(\varepsilon^2)\)，\(S=R/(\varepsilon)\)。模 \(S\) 置于零次虽然有界且有限生成，却不是完美复形。它的自由消解为
\[
 \cdots\xrightarrow{\varepsilon}R\xrightarrow{\varepsilon}R
 \xrightarrow{\varepsilon}R\longrightarrow S\longrightarrow0,
\]
因为每个乘 \(\varepsilon\) 映射的核、像都是 \(\varepsilon R\)。施加 \(\operatorname{Hom}_R(-,S)\) 后微分全为零，所以 \(\operatorname{Ext}^n_R(S,S)=k\) 对所有 \(n\ge0\) 成立。若 \(S\) 有有界投射模型 \(P\)，则 \(\underline{\operatorname{Hom}}_R(P,S)\) 有界，其高次上同调终将消失；结合\cref{b4:thm:derived-ext-identification}即得矛盾。
\end{example}

\subsection{厚子范畴}
\label{b4:subsec:2002:02}

在有界有限生成投射复形中，有限直和、移位和锥并不能单独表达“取一个投射直和项”的操作。例如，从自由模 \(R^n\) 取出有限生成投射直和项，是完美对象形成过程的一部分，不能只允许取自由项。

\begin{definition}\label{b4:def:thick-subcategory}
设 \(\mathcal T\) 为三角范畴。一个含零对象、对同构、移位、逆移位和映射锥封闭的全子范畴 \(\mathcal S\)，若还满足：当 \(X\in\mathcal S\) 且 \(X\cong Y\oplus Z\) 在 \(\mathcal T\) 中成立时，\(Y,Z\in\mathcal S\)，便称为\term[houzifanchou]{厚子范畴}[thick subcategory]。包含对象族 \(\mathcal U\) 的最小厚子范畴记为 \(\operatorname{thick}(\mathcal U)\)。
\end{definition}
\symidx{\(\operatorname{thick}(T)\)：由对象生成的厚子范畴}

厚生成只允许有限次构造以及有限直和；没有规定可以取无限直和。“某对象属于 \(\operatorname{thick}(T)\)”因而是一条有限构造声明。

\begin{proposition}\label{b4:prop:perfect-built-from-ring}
设 \(R\) 为含幺环，\(P\) 为有界有限生成投射左 \(R\)-模复形。在 \(D(R\text{-}\mathrm{Mod})\) 中有 \(P\in\operatorname{thick}(R)\)，这里 \(R\) 作为集中在零次的左正则模。
\end{proposition}
\begin{proof}
每个 \(P^n\) 都是某个有限自由模 \(R^{r_n}\) 的直和项，故 \(P^n[-n]\in\operatorname{thick}(R)\)。按非零项的数目归纳。若最低非零次数为 \(a\)，取从 \(a+1\) 次开始保留原复形的子复形 \(P^{\ge a+1}_{\mathrm{br}}\)。有逐次分裂短正合列
\[
 0\longrightarrow P^{\ge a+1}_{\mathrm{br}}
 \longrightarrow P\longrightarrow P^a[-a]\longrightarrow0.
\]
它在同伦范畴、因而也在导出范畴中给出好三角。归纳假设处理左项，右项已属于所需厚子范畴，于是中项也属于其中。
\end{proof}

反向包含有一个实质细节：导出范畴中的直和分解，最初只给出同伦意义的幂等态射，未必已经是所选模型逐项分裂的幂等矩阵。下一小节的证明概要说明，怎样从紧性提取有限构造，再把这个直和项写成有界有限生成投射复形。

\subsection{紧性概念的适用背景}
\label{b4:subsec:2002:03}

紧性比较一个对象到任意直和的态射与它到各直和项的态射。因此必须先有任意集合指标的余积。这里采用 \(D(R\text{-}\mathrm{Mod})\)，不把讨论含混地放在可能没有无限直和的有界有限生成模范畴里。

\begin{definition}\label{b4:def:compact-object}
设 \(\mathcal T\) 为具有任意集合指标余积的局部小三角范畴，\(P\in\mathcal T\)。若对每族对象 \((X_i)_{i\in I}\)，自然映射
\[
 \bigoplus_{i\in I}\operatorname{Hom}_{\mathcal T}(P,X_i)
 \longrightarrow\operatorname{Hom}_{\mathcal T}\left(P,\bigoplus_{i\in I}X_i\right)
\]
为同构，则称 \(P\) 为\term[jinduixiang]{紧对象}[compact object]。若对象 \(G\in\mathcal T\) 满足：对任意 \(X\in\mathcal T\)，\(\operatorname{Hom}_{\mathcal T}(G,X[n])=0\) 对所有 \(n\in\mathbb Z\) 成立时必有 \(X=0\)，则称 \(G\) 为检测对象是否为零的生成元。
\end{definition}

模范畴的直和正合，因而复形的逐次直和保持拟同构，并给出模导出范畴的余积。后一断言还须检验态射的泛性质：对一族 \(X_i\)，用\cref{b4:thm:semifree-resolution}取半自由替换 \(P_i\to X_i\)。直和 \(P=\bigoplus_iP_i\) 仍为半自由复形，且 \(P\to\bigoplus_iX_i\) 拟同构。由\cref{b4:prop:k-resolution-comparison}，从 \(P\) 出发的导出态射可以在同伦范畴计算；而从直和出发的复形映射及其同伦都是逐分量给定的，故
\[
 \operatorname{Hom}_D\left(\bigoplus_iX_i,Y\right)
 \cong\prod_i\operatorname{Hom}_K(P_i,Y)
 \cong\prod_i\operatorname{Hom}_D(X_i,Y).
\]
这就是余积泛性质。DG 模的相应构造可对照\cite[Sections 3.1 and 4.1]{Keller1994DerivingDG}；两种论证都不能以任意逐项投射复形替代所需的无界替换。

\begin{proposition}\label{b4:prop:perfect-compact-ring-generator}
设 \(R\) 为含幺环。在 \(D(R\text{-}\mathrm{Mod})\) 中，每个完美复形都是紧对象。左正则模 \(R\) 还是检测对象是否为零的生成元。
\end{proposition}
\begin{proof}
取有界有限生成投射模型 \(P\)，令 \(S=\{\,n\in\mathbb Z\mid P^n\ne0\,\}\)，这是有限集合。固定 Hom 次数 \(r\)，有
\[
\begin{aligned}
 \underline{\operatorname{Hom}}_R\left(P,\bigoplus_iX_i\right)^r
 &=\prod_{n\in S}\operatorname{Hom}_R
       \left(P^n,\bigoplus_iX_i^{n+r}\right)\\
 &\cong\prod_{n\in S}\bigoplus_i\operatorname{Hom}_R(P^n,X_i^{n+r})\\
 &\cong\bigoplus_i\prod_{n\in S}\operatorname{Hom}_R(P^n,X_i^{n+r}).
\end{aligned}
\]
第一处同构使用 \(P^n\) 有限生成：有限个生成元的像总共只涉及有限个直和项。第二处使用 \(S\) 有限，使各次数所涉及的有限指标集可以合成一个有限集。这正是有界性在 Hom 总化中的作用；若直积指标无限，不能如此交换直积与直和。上述同构与 Hom 微分相容，因而自然映射
\[
 \bigoplus_i\underline{\operatorname{Hom}}_R(P,X_i)
 \longrightarrow
 \underline{\operatorname{Hom}}_R\left(P,\bigoplus_iX_i\right)
\]
是复形同构，对它取 \(H^0\) 仍得到同构。左边逐分量取圈与边界，便得到各 Hom 复形零次上同调的直和；有界投射模型的全忠实性把上述同构识别为紧性的映射。另一方面，\(\operatorname{Hom}_{D(R\text{-}\mathrm{Mod})}(R,X[n])=H^n(X)\)。若全部为零，\(X\) 零调，便在导出范畴中同构于零。
\end{proof}

\begin{theorem}[完美复形的内在刻画]\label{b4:thm:perfect-compact-thick}
设 \(R\) 为含幺环。在 \(D(R\text{-}\mathrm{Mod})\) 中，下列三类对象相同：完美复形、紧对象以及 \(\operatorname{thick}(R)\) 中的对象。特别地，完美复形组成厚子范畴。
\end{theorem}

\begin{proofsketch}
完美推出紧、完美属于 \(\operatorname{thick}(R)\) 已分别由\cref{b4:prop:perfect-compact-ring-generator,b4:prop:perfect-built-from-ring}证明。紧对象对移位、有限直和、锥和直和项封闭：对锥使用 Hom 长正合列和五引理，对直和项限制紧性的自然同构即可。因此 \(\operatorname{thick}(R)\) 中的对象都是紧对象。只剩紧推出完美。

设 \(C\) 紧，用\cref{b4:thm:semifree-resolution}取半自由消解 \(P\to C\)，其滤过 \(F_0P\subseteq F_1P\subseteq\cdots\) 的各次商都是若干移位 \(R[n]\) 的直和。下列望远镜三角表示 \(P\)：
\[
 \bigoplus_{j\ge0}F_jP\xrightarrow{\,1-s\,}
 \bigoplus_{j\ge0}F_jP\longrightarrow P
 \longrightarrow\left(\bigoplus_{j\ge0}F_jP\right)[1],
\]
其中 \(s\) 把第 \(j\) 项经滤过包含送到第 \(j+1\) 项。在每个复形次数上，\((1-s)(x)=0\) 的零号坐标给出 \(x_0=0\)，随后递推得到 \(x_1=x_2=\cdots=0\)，所以 \(1-s\) 单射。其余核正是按滤过包含取的余极限，即 \(P\)，因而这一逐次短正合列给出所写好三角。对它取 \(\operatorname{Hom}(C,-)\)，紧性把两端的 Hom 变成直和；相同的坐标递推说明，后一移位的 \(1-s\) 仍单射。Hom 长正合列遂得
\[
 \operatorname*{colim}_j\operatorname{Hom}_D(C,F_jP)
 \cong\operatorname{Hom}_D(C,P).
\]
令 \(u:C\xrightarrow{\sim}P\) 为消解拟同构在导出范畴中的逆。上述余极限公式给出某个有限阶段及 \(v:C\to F_jP\)，使自然映射 \(a:F_jP\to P\) 满足 \(av=u\)。于是
\[
 (u^{-1}a)v=1_C.
\]
这说明 \(C\) 是 \(F_jP\) 的导出直和项。这里仅有滤过步数有限，每一层的自由生成元仍可能有无限多个。

还需把各层的无限直和缩为有限直和，并保留连接态射。两层时可以把这一步具体写出。设
\[
 U\xrightarrow{i}E\xrightarrow{p}V\xrightarrow{\delta}U[1]
\]
为好三角，其中 \(U,V\) 均为移位 \(R[n]\) 的任意直和，给定 \(f:C\to E\)。紧性使 \(pf\) 经某个有限子直和 \(V_0\) 分解为 \(pf=j_Vg\)，其中 \(j_V:V_0\to V\) 为包含。\(V_0\) 本身完美、因而紧，所以 \(\delta j_V\) 又经 \(U[1]\) 中某个有限子直和 \(U_0[1]\) 分解：
\[
 \delta j_V=j_U[1]\delta_0,
 \qquad \delta_0:V_0\longrightarrow U_0[1].
\]
这里 \(j_U:U_0\to U\) 是自由直和的分裂包含。由 \(\delta pf=0\) 得 \(j_U[1]\delta_0g=0\)，再用其左逆才得到 \(\delta_0g=0\)。把 \(\delta_0\) 补成好三角
\[
 U_0\longrightarrow E_0\xrightarrow{p_0}V_0
 \xrightarrow{\delta_0}U_0[1].
\]
\(E_0\) 由有限个 \(R[n]\) 用一次锥构造得到。态射公理给出 \(q:E_0\to E\)，与两端的包含相容；\cref{b4:prop:triangle-hom-exact}给出 \(v_0:C\to E_0\)，使 \(p_0v_0=g\)。因此 \(p(f-qv_0)=0\)，再次用 Hom 正合性得到 \(t:C\to U\)，使 \(f-qv_0=it\)。最后用 \(C\) 紧，将 \(t\) 写成 \(t=j_1t_1\)，其中 \(j_1:U_1\to U\) 是另一个有限子直和的包含。这样
\[
 f=\begin{pmatrix}q&ij_1\end{pmatrix}
   \begin{pmatrix}v_0\\t_1\end{pmatrix}
 :C\longrightarrow E_0\oplus U_1\longrightarrow E
\]
就是所需的有限因子分解。除了压缩两端，还须控制提升的障碍 \(\delta_0g\)，以及提升以后与原映射之差；这两处不能只靠“取有限子直和”略去。

一般有限层数使用 Keller \cite[Section 5.3, Lemma]{Keller1994DerivingDG}所给的有限因子分解引理及其八面体归纳。该引理同时处理态射与锥：若 \(c:Z'\to Z\) 的锥由自由模移位的任意直和经有限次扩张得到，则从紧对象出发的 \(C\to Z\) 可补成从 \(b:C'\to C\) 到 \(c\) 的交换方块，使 \(b\) 的锥只由有限个自由模移位经有限次扩张得到。应用于 \(0\to F_jP\) 及上述 \(v\)，交换性给出 \(vb=0\)；再由 \(v\) 有左逆得 \(b=0\)。\cref{b4:prop:split-triangle-criterion}遂使 \(C\) 成为 \(\operatorname{Cone}(b)\) 的直和项，故 \(C\in\operatorname{thick}(R)\)。至此，已得到紧对象恰为 \(\operatorname{thick}(R)\) 中的对象；还没有把这些导出直和项写成有界有限生成投射模型。

要得到完美模型，还需处理同伦意义的幂等态射。若 \(e:R^m\to R^m\) 已是满足 \(e^2=e\) 的模同态，则 \(R^m=\operatorname{im}e\oplus\ker e\)，其像直接是有限生成投射模。可是在有界投射复形 \(Q\) 的模型上，导出直和项只给出 \([e]^2=[e]\)，即 \(e^2-e\) 零同伦；这不保证每个次数分量 \(e^n:Q^n\to Q^n\) 满足 \((e^n)^2=e^n\)，不能直接逐次取其像作为所需复形。

这里引用 Keller 对负 DG 范畴的加强结论在普通环上的版本：紧对象本身可以写成有限次扩张，且各扩张因子是有限个自由模移位之和的导出直和项\cite[Section 5.3, Remark b]{Keller1994DerivingDG}。原处说明此加强结论由 Rickard 的证明改编得到；普通环的紧对象与有界有限生成投射模型的对应也见\cite[Section 3.1]{Keller1998AbelianDerived}。普通环集中在零次，满足正次数态射复形为零的负 DG 条件。这里先将对象重新组织为有限扩张，再对每个扩张因子取直和项；不能从 \(C\) 是某个有限扩张的直和项，直接推断总对象的幂等态射保持原来的扩张滤过。

对上述任一因子 \(G\)，把它写成 \(V=\bigoplus_{j=1}^m R[n_j]\) 的导出直和项。不同移位之间没有非零态射，即
\[
\operatorname{Hom}_D(R[n],R[n'])=0\qquad(n\ne n').
\]
因此 \(V\) 上对应的幂等态射按相同移位分块。在每一块上，它由有限自由左 \(R\)-模的真正幂等自同态表示，其像是有限生成投射模。因此 \(G\) 同构于有限个有限生成投射模移位的直和。最后，有界投射模型的全忠实性使各次扩张的连接态射都能由复形映射表示，逐次取锥便得到 \(C\) 的有界有限生成投射模型。

以上概要省略了有限因子分解引理的一般八面体归纳，并引用了处理导出直和项的有限扩张结论。前者及紧对象论证见\cite[Sections 5.1--5.3, especially the lemma in 5.3]{Keller1994DerivingDG}；后者用来完成最后一项 \(\operatorname{thick}(R)\subseteq\operatorname{Perf}(R)\)，并非已经由逐项的投射模分解证明。
\end{proofsketch}

两个例子说明条件各有作用。若 \(R\ne0\)，无限自由模 \(F=\bigoplus_{j\ge0}R\) 不是紧对象：恒等映射 \(F\to\bigoplus_{j\ge0}R\) 不经过任何有限子直和。另一方面，\cref{b4:exm:bounded-not-perfect}中的单模有限生成且集中在一个次数，但其无限投射消解妨碍紧性。只检查“每项有限”或者“同调有界”，都不能替代完美性。
